\section{Arbitrary kernels on a five-point Sidon support}
\label{sec:five-point-weighted}

The support that defeats the universal covering bound also satisfies the
reflected inequality for every choice of kernel weights. The proof uses
rational coefficient certificates in addition to the estimates below.

\begin{theorem}\label{thm:five-point-weighted}
Let $F=\{0,2,7,8,11\}$. Every nonnegative kernel $k$ supported on $F$
and every pair of nonnegative finitely supported functions $f,g$ on
$\mathbb Z$ satisfy
\[
 \|f\star\widetilde{gk}\|_1^4
 \le T_F(k)\,\|f\star g\|_1^4.
\]
The inequality is strict when $k$ is positive at all five points and
$f,g$ are nonzero.
\end{theorem}

This result concerns the displayed support. Theorem~\ref{thm:five-point}
allows every five-point integer support when $k$ is constant. Neither
theorem assumes a bound on the size or diameter of the support of $f$.

\subsection{Allocations from a common input}

Initially suppose that $f\ne0$ and that $g$ is positive on a five-point
set $F=\{p_1,\ldots,p_5\}$ and zero elsewhere. Write
$P=\|f\star g\|_1>0$. For each positive term of
\[
 N=\|f\star\widetilde{gk}\|_1
   =\sum_y\max_i k_i g_i f(y+p_i),
\]
select a maximizing index, resolving ties by index order. Let $E_i$
be the disjoint sets of outputs assigned to index $i$, and put
\[
 a_i=\frac{g_i}{P}\sum_{y\in E_i}f(y+p_i).
\]
Then $N/P=\sum_i k_i a_i$. Denominator outputs at $y+2p_i$ give
$0\le a_i\le1$. For distinct $i,j$, translate both $E_i$ and $E_j$
by $p_i+p_j$. They remain disjoint, and the corresponding denominator
witnesses give
\begin{equation}\label{weighted-pairs}
 \frac{g_j}{g_i}a_i+\frac{g_i}{g_j}a_j\le1,
 \qquad a_i a_j\le\frac14.
\end{equation}
Thus these constraints hold simultaneously for allocations obtained
from the same $f,g$, even though their selection depends on $k$.

\begin{lemma}\label{weighted-total}
Every allocation above satisfies
\[
 \sum_i a_i\le c:=\frac{1109}{450}.
\]
\end{lemma}
\begin{proof}
Put $R=\max_i g_i/\min_i g_i$ and $t=1/R$. We first show that
\begin{equation}\label{weighted-range}
 A:=\sum_i a_i\le F(t):=2+t-\frac{t^2}{2}.
\end{equation}
The pair at the largest and smallest values of $g$ satisfies
$ta+b/t\le1$. If $a,b\le v\le1$, filling the coordinate with
the smaller coefficient first gives
\[
 a+b\le\min\{2v,v+t-vt^2\}.
\]
When $M:=\max_i a_i\le1/2$, take $v=1/2$ and bound the other
three coordinates by $1/2$ to obtain \eqref{weighted-range}.

For $1/2\le M\le1$, all other coordinates are at most $1/(4M)$
by \eqref{weighted-pairs}. If the maximum is outside the extremal
pair, the preceding two-variable bound gives
\[
 A\le\min\{U(M),V(M)\},\qquad
 U(M)=M+\frac1M,\quad V(M)=M+\frac{3-t^2}{4M}+t.
\]
If the maximum belongs to the pair, assigning it to the smaller
coefficient can only increase the allowed partner. This gives the
same first bound and the second bound
$M+3/(4M)+t-Mt^2\le V(M)$.

The functions switch order at $M_0=(1+t^2)/(4t)\ge1/2$.
The function $U$ decreases on $[1/2,1]$, while $V$ is convex there.
Consequently it suffices to check $V(1/2)=F(t)$,
$V(1)\le F(t)$, and, if $M_0\le1$, $U(M_0)\le F(t)$.
The endpoint difference is $F(t)-V(1)=(1-t^2)/4$. At the switch,
\[
 4t(1+t^2)\bigl(F(t)-U(M_0)\bigr)
 =(1-t)^2(-2t^3-t^2+6t-1).
\]
Here $M_0\le1$ implies $t\ge2-\sqrt3>1/4$. On $[1/4,1]$,
the last factor is at least
$-3t^2+6t-1=2-3(1-t)^2\ge5/16$. This proves
\eqref{weighted-range}.

At each output, the selected unweighted product is at most
$\max_i g_i f(y+p_i)$. Therefore the energy estimate
\eqref{five-range-energy} also gives $A\le9R/5$. Combining the
two bounds yields
\[
 A\le\min\left\{\frac{9R}{5},\,
                  2+\frac1R-\frac1{2R^2}\right\}.
\]
For $R\le15/11$, the first expression is at most $27/11<c$.
For $R\ge15/11$, the second is decreasing and is at most $c$.
\end{proof}

For the kernel, write $S=\sum_i k_i$, $K=\max_i k_i$, and
$m=\min_i k_i$. Define
\[
 \delta=\frac52-c=\frac8{225},\qquad
 \alpha=\frac{c-\sqrt{c^2-4}}2,\qquad
 \beta=\frac1{4\alpha}.
\]
Then $1/2<\alpha<1$ and $\alpha+4\beta=c$.

\begin{lemma}\label{weighted-envelope}
The allocations satisfy
\[
 \sum_i k_i a_i\le\max\{U_c,V_c,C\},\qquad
 U_c=\frac S2-\delta m,\quad
 V_c=\alpha K+\beta(S-K),\quad C=\frac{S+3K}{4}.
\]
\end{lemma}
\begin{proof}
Use only $0\le a_i\le1$, $a_i a_j\le1/4$, and $\sum_i a_i\le c$.
These constraints are invariant under coordinate permutations, so
rearrangement permits a largest allocation $M$ at a largest kernel weight.
If $M\le1/2$, the deficit from five halves is at least $\delta$;
its weighted cost is at least $\delta m$, giving $U_c$.

If $1/2\le M\le\alpha$, the other four coordinates are at most
$1/(4M)$. Since $M+1/M\ge c$, subtracting the least possible cost
of the missing mass bounds the objective by
\[
 mc+(K-m)M+\frac{S-K-4m}{4M}.
\]
The coefficient $S-K-4m$ is nonnegative, so this expression is convex
in $M$. Its endpoint values are $U_c$ and $V_c$.
For $\alpha\le M\le1$, the convex bound
$KM+(S-K)/(4M)$ has endpoint values $V_c$ and $C$.
\end{proof}

\subsection{The corner comparison}

\begin{lemma}\label{weighted-corner}
If $F$ is a five-point Sidon set with $|4F|\ge32$, then
\[
 T_F(k)\ge\left(\frac{S+3K}{4}\right)^4.
\]
The inequality is strict when at least two kernel coordinates are positive.
\end{lemma}
\begin{proof}
The zero kernel is immediate. Normalize $K=1$, choose a position $p_i$
with weight one, and let $u_1,\ldots,u_4$ be the other weights.
Set $t=\sum_j u_j\in[0,4]$. Sidonicity makes the 15 outputs
$2p_i+(p_j+p_\ell)$, $j\le\ell$, distinct. Their total contribution
is at least
\[
 1+t+\frac{t^2+\sum_j u_j^2}{2}\ge1+t+\frac{5t^2}{8}.
\]
At least 17 other fourfold outputs each contribute at least $m^4$,
where $m\ge(t-3)_+$. Hence
\[
 T_F(k)-(1+t/4)^4
 \ge t^2 A(t)+17(t-3)_+^4,\qquad
 A(t)=\frac14-\frac t{16}-\frac{t^2}{256}.
\]
The function $A$ decreases on $[0,4]$. For $0<t\le13/4$ it is
at least $23/4096>0$. On the remaining intervals, the displayed
lower bound is bounded below respectively by
\[
\begin{array}{c|c}
 [13/4,10/3]&-25/1296+17/256=977/20736\\[2pt]
 [10/3,7/2]&-833/4096+17/81=2159/331776\\[2pt]
 [7/2,4]&-1+17/16=1/16.
\end{array}
\]
All are positive. When $t=0$, both sides equal one.
\end{proof}

For $F=\{0,2,7,8,11\}$, all 15 unordered pair sums are distinct,
and direct addition gives
\begin{equation}\label{weighted-fourfold-support}
 4F=\{0,2,4\}\cup\{6,7,\ldots,38\}\cup\{40,41,44\}.
\end{equation}
Thus $|4F|=39$ and Lemma~\ref{weighted-corner} applies.

\subsection{Rational certificates on the weight orders}

It remains to compare both $U_c$ and $V_c$ with $T_F(k)^{1/4}$.
Use the rational majorant
\[
 \overline V=aK+b(S-K),\qquad
 a=\frac{41}{80},\quad b=\frac{7027}{14400}=\frac{c-a}{4}.
\]
Indeed $a+1/a=8081/3280<c$, so $\alpha<a$. The function
$uK+(c-u)(S-K)/4$ has derivative $(5K-S)/4\ge0$;
therefore $V_c\le\overline V$.

Index the five positions of $F$ in increasing order. On a cone where
their weights have order $\pi$, write
\[
 k_{\pi(j)}=m+d_1+\cdots+d_{j-1}\quad(1\le j\le5),
 \qquad x=(m,d_1,d_2,d_3,d_4)\ge0.
\]
There are $5!=120$ orders, with ties allowed. The two linear targets
$U_c=q^U\cdot x$ and $\overline V=q^V\cdot x$ have vectors
\begin{align*}
 q^U&=(1109,900,675,450,225)/450,\\
 q^V&=(35488,28461,21434,14407,7380)/14400.
\end{align*}

Let $\mathcal M=\{\mu\in\mathbb N_0^5:|\mu|=4\}$ be the 70
unordered fourfold multisets. For $\mu\in\mathcal M$, put
\[
 z(\mu)=\sum_i p_i\mu_i,\qquad W_\mu(k)=\prod_i k_i^{\mu_i},
 \qquad \mathcal M_z=\{\mu:z(\mu)=z\}.
\]
There are 39 nonempty fibers $\mathcal M_z$ by
\eqref{weighted-fourfold-support}. These witness monomials have no
ordered-representation multiplicity: the operation is max-convolution.

For each order $\pi$ and each target $q$, the certificate supplies 70
nonnegative integers $n_\mu$, with common denominator $D=10^9$,
satisfying
\begin{equation}\label{weighted-certificate}
 \sum_{\mu\in\mathcal M_z}n_\mu=D\quad(z\in4F),\qquad
 \sum_\mu\frac{n_\mu}{D}W_\mu(k(x))-(q\cdot x)^4
 \text{ has nonnegative coefficients.}
\end{equation}
The supplied verifier checks these conditions for all 240 pairs
$(\pi,q)$. Its arithmetic uses integers throughout each coefficient
comparison. If $q=v/Q$, the comparison for a degree-four exponent
$\gamma$ is
\[
 Q^4\sum_\mu n_\mu[x^\gamma]W_\mu(k(x))
       -D[x^\gamma](v\cdot x)^4\ge0.
\]
The program reconstructs every witness coefficient by enumerating the
variable choices in its four slots, and expands the target by repeated
polynomial multiplication. It checks the fiber masses, all 70 coefficients,
and both targets for every order. The coefficient inequalities are checked
independently of the solver's floating-point values.

For every $x\ge0$, \eqref{weighted-certificate} implies
\[
 T_F(k)=\sum_z\max_{\mu\in\mathcal M_z}W_\mu(k)
 \ge\sum_\mu\frac{n_\mu}{D}W_\mu(k)
 \ge(q\cdot x)^4.
\]
Every witness has $m^4$ coefficient one. Fiber normalization therefore
makes that coefficient 39 in the middle polynomial, whereas both target
fourth powers have coefficient $c^4<3721/100<39$. Thus both comparisons
are strict when $m>0$.

\begin{proof}[Proof of Theorem~\ref{thm:five-point-weighted}]
For positive $k,g$ on $F$ and nonzero $f$, combine the allocation
envelope with the two certificate comparisons and the corner lemma.
Each of $U_c,V_c,C$ is strictly below $T_F(k)^{1/4}$, so the
reflected inequality is strict.

For fixed $f$, all relevant norms are finite sums of maxima of finitely
many coordinate products, and hence continuous in $k,g$. Positive
approximation gives the non-strict inequality when some coordinates
vanish. When $k>0$, the upper envelope depends only on $k$ and stays
strictly below $T_F(k)^{1/4}$; taking the same limit in the envelope
therefore preserves strictness for nonzero $g$ supported on $F$.

For general $g$, restrict it to $g_F=g1_F$. The numerator is unchanged,
while $\|f\star g_F\|_1\le\|f\star g\|_1$. If $g_F\ne0$, apply
the preceding result. If $g_F=0$, the numerator is zero; when $k>0$
and $f,g\ne0$, the right side is positive. The zero-function cases
of the non-strict statement are immediate.
\end{proof}

The certificate files and standard-library checker accompany the paper.
From this paper's source directory, run
\begin{verbatim}
python3 -B verify_five_point.py
\end{verbatim}
The command checks all 240 certificates and corruption controls. It
requires no numerical optimizer or external Python package. The finite
certificates establish the coefficient comparisons; the preceding hand
argument supplies their connection to the reflected inequality.
